% TOP_PROBLEM: {"id":30007744,"problem_number":"LOCAL-30007744","title":"Construction-productivity stagnation","category_id":21,"category":{"id":21,"name":"economics","display_name":"Economics","description":"Open economic theory statements, published research agendas, and proposed scoped specifications; question provenance and historical priority are qualified per record.","slug":"economics","order_index":21,"created_at":"2026-10-08T00:00:00Z"},"status":"research_question","external_url":null,"legacy_ids":[],"published":false,"statement":"Identify how a defined land-use streamlining reform changes construction productivity while separating quality measurement, project scale, technology, and firm reallocation.","economics":{"original_id":233,"field":"Urban, housing and spatial economics","kind":"identification","formalization_basis":"adapted_research_question","source_status":"research_agenda","priority_status":"earliest_found","priority_note":"Earliest explicit statement located in this audit; no exhaustive priority search or first-ever claim. The mathematical specification below is adapted, not quoted from the source.","earliest_verified_reference":{"authors":["Nathaniel Baum-Snow","Gilles Duranton"],"title":"Housing supply and housing affordability","year":2025,"url":"https://realestate.wharton.upenn.edu/wp-content/uploads/2025/12/BaumSnow_Duranton_bc_2025-003.pdf","doi":"10.1016/bs.hesreg.2025.05.003","locator":"§7 Conclusions and future research directions, printed pp. 448–449","evidence_summary":"The authors explicitly seek dynamic supply elasticities, construction-productivity causes and remedies, and richer spatial models of migration and housing reform."},"current_reference":{"authors":["Nathaniel Baum-Snow","Gilles Duranton"],"title":"Housing supply and housing affordability","year":2025,"url":"https://realestate.wharton.upenn.edu/wp-content/uploads/2025/12/BaumSnow_Duranton_bc_2025-003.pdf","doi":"10.1016/bs.hesreg.2025.05.003","locator":"§7 Conclusions and future research directions, printed pp. 448–449","evidence_summary":"The authors explicitly seek dynamic supply elasticities, construction-productivity causes and remedies, and richer spatial models of migration and housing reform."},"formalization_note":"The source explicitly identifies the underlying gap or agenda. Population, horizon, interventions, estimands, and auxiliary assumptions are proposed operational choices, not a canonical source theorem. Partial later answers are compatible with this scoped question; the citation alone does not prove absence of every subsequent solution. The 2024 regulation paper supplies an explanation, while the explicit general productivity research call inspected here is in the 2025 housing survey. Mathematical scope was checked independently; added definedness, feasibility, or identification conditions belong to this proposed formalization. Mathematical scope was checked independently; added definedness, feasibility, or identification conditions belong to this proposed formalization. Mathematical scope was checked independently; added definedness, feasibility, or identification conditions belong to this proposed formalization. Mathematical scope was checked independently; added definedness, feasibility, or identification conditions belong to this proposed formalization."},"statusQualification":"Published question or research agenda verified at the cited locator. The mathematical specification is adapted for this catalogue and includes additional assumptions; source publication does not establish first-ever priority or certify this exact model remains unsolved. The source explicitly identifies the underlying gap or agenda. Population, horizon, interventions, estimands, and auxiliary assumptions are proposed operational choices, not a canonical source theorem. Partial later answers are compatible with this scoped question; the citation alone does not prove absence of every subsequent solution. The 2024 regulation paper supplies an explanation, while the explicit general productivity research call inspected here is in the 2025 housing survey. Mathematical scope was checked independently; added definedness, feasibility, or identification conditions belong to this proposed formalization. Mathematical scope was checked independently; added definedness, feasibility, or identification conditions belong to this proposed formalization. Mathematical scope was checked independently; added definedness, feasibility, or identification conditions belong to this proposed formalization. Mathematical scope was checked independently; added definedness, feasibility, or identification conditions belong to this proposed formalization.","statusReviewedAt":"2026-10-08"}
% FIRST_POSTED: 2026-10-08
% ENTRY_KIND: statement_only
% !TeX program = lualatex
\documentclass[11pt]{article}
\usepackage{fontspec}
\usepackage[a4paper,margin=25mm]{geometry}
\usepackage{amsmath,amssymb,mathtools}
\usepackage{xurl}
\usepackage[hidelinks]{hyperref}
\newcommand{\sref}[2]{\hyperlink{source:#1:#2}{[S#2]}}
\setlength{\emergencystretch}{3em}
\begin{document}
% problemId:30007744
\section*{Construction-productivity stagnation}\label{30007744}
\subsection{Definitions and mathematical statement}
Consider residential construction establishments in one country from 2000 through 2025. Let \(Q_{it}\) be quality-adjusted completed floor area, \(L_{it}\) labor hours, \(K_{it}\) capital services, and \(M_{it}\) materials for establishment \(i\) in year \(t\). For operating establishments with positive output and input aggregator, define productivity \(A_{it}=Q_{it}/F(L_{it},K_{it},M_{it})\), where the production aggregator \(F\) and quality adjustments are stated and estimated independently of the policy contrast. Let \(r\in\{0,1\}\) denote current land-use rules versus a specified streamlined regime. Let \(O_i(r)\) indicate operation at follow-up under regime \(r\). The survivor causal target is
\[\tau_A=E[\log A_{i,t+5}(1)-\log A_{i,t+5}(0)\mid O_i(1)=O_i(0)=1].\]
Assume the always-operating stratum has positive probability and integrable log contrasts. Its membership is latent: use justified principal-stratum assumptions or bounds, and report survival effects separately. Identify whether regulation changes productivity through project scale, establishment entry, standardization, or technology adoption, rather than merely changing output prices or building composition. Track subcontracting and unrecorded hours; aggregate industry productivity additionally requires weights and an entry treatment. Use credible rule changes with justified counterfactual trends and test sensitivity to alternative output deflators. Decompose aggregate stagnation only when within-firm changes, reallocation, and quality are separately identified. The housing survey explicitly requests research on construction-productivity causes and remedies. The source does not provide this exact establishment-level estimand; its target and measurement restrictions are an adapted specification.

\par\textbf{Formulation and scope.} The source explicitly identifies the underlying gap or agenda. Population, horizon, interventions, estimands, and auxiliary assumptions are proposed operational choices, not a canonical source theorem. Partial later answers are compatible with this scoped question; the citation alone does not prove absence of every subsequent solution. The 2024 regulation paper supplies an explanation, while the explicit general productivity research call inspected here is in the 2025 housing survey. Mathematical scope was checked independently; added definedness, feasibility, or identification conditions belong to this proposed formalization. Mathematical scope was checked independently; added definedness, feasibility, or identification conditions belong to this proposed formalization. Mathematical scope was checked independently; added definedness, feasibility, or identification conditions belong to this proposed formalization. Mathematical scope was checked independently; added definedness, feasibility, or identification conditions belong to this proposed formalization.

\par\textbf{Open-question provenance.} An explicit question or research agenda was verified at the source locator. This does not by itself certify that every later version remains unresolved \sref{30007744}{1}.

\par\textbf{Historical priority.} Earliest explicit statement located in this audit; no exhaustive priority search or first-ever claim. The mathematical specification below is adapted, not quoted from the source.

\subsection{Short English statement}
Identify how a defined land-use streamlining reform changes construction productivity while separating quality measurement, project scale, technology, and firm reallocation.

\subsection{Sources}
\begin{itemize}\raggedright
\item[\textnormal{[S1]}]\hypertarget{source:30007744:1}{} Nathaniel Baum-Snow, Gilles Duranton. 2025. Housing supply and housing affordability. §7 Conclusions and future research directions, printed pp. 448–449.
\url{https://realestate.wharton.upenn.edu/wp-content/uploads/2025/12/BaumSnow_Duranton_bc_2025-003.pdf}.
DOI: 10.1016/bs.hesreg.2025.05.003.
{\small\textit{Earliest verified question/agenda reference; historical priority qualified below. The authors explicitly seek dynamic supply elasticities, construction-productivity causes and remedies, and richer spatial models of migration and housing reform.}}
\end{itemize}
\end{document}
